\documentclass[10pt,t]{beamer}
\usepackage[fontset=ubuntu]{ctex}
\usepackage{graphicx}
\usepackage{listings}
\usepackage{ulem}
\usepackage{boxedminipage}
% Copyright 2007 by Till Tantau
%
% This file may be distributed and/or modified
%
% 1. under the LaTeX Project Public License and/or
% 2. under the GNU Public License.
%
% See the file doc/licenses/LICENSE for more details.


% Common packages

%\usepackage{babel}
%\usepackage[latin1]{inputenc}
\usepackage{times}
\mode<article>
{
  \usepackage{times}
  \usepackage{mathptmx}
  \usepackage[left=1.5cm,right=6cm,top=1.5cm,bottom=3cm]{geometry}
}

\usepackage{hyperref}
%\usepackage[T1]{fontenc}
%\usepackage{tikz}
%\usepackage{colortbl}
%\usepackage{translator} % comment this, if not available


% Common settings for all lectures in this course

\def\lecturename{}

\title{\insertlecture}

\author{}

\institute
{
    Fang Yuan (yfang@nju.edu.cn)\\
  School of Electronics Science and Engineering\\
  Nanjing University
}

\subject{\lecturename}


% Beamer version theme settings

\useoutertheme[height=0pt,width=2cm,right]{sidebar}
\usecolortheme{rose,sidebartab}
\useinnertheme{circles}
\usefonttheme[only large]{structurebold}

\setbeamercolor{sidebar right}{bg=black!15}
\setbeamercolor{structure}{fg=blue}
\setbeamercolor{author}{parent=structure}

\setbeamerfont{title}{series=\normalfont,size=\LARGE}
\setbeamerfont{title in sidebar}{series=\bfseries}
\setbeamerfont{author in sidebar}{series=\bfseries}
\setbeamerfont*{item}{series=}
\setbeamerfont{frametitle}{size=}
\setbeamerfont{block title}{size=\small}
\setbeamerfont{subtitle}{size=\normalsize,series=\normalfont}

\setbeamertemplate{navigation symbols}{}
\setbeamertemplate{bibliography item}[book]
\setbeamertemplate{sidebar right}
{
  {\usebeamerfont{title in sidebar}%
%    \vskip1.5em%
    \hskip3pt%
    \usebeamercolor[fg]{title in sidebar}%
%    \insertshorttitle[width=2cm,center,nolinebreaks]\par%
    \vskip1.25em%
  }%
  {%
    \hskip3pt%
    \usebeamercolor[fg]{author in sidebar}%
    \usebeamerfont{author in sidebar}%
    \insertshortauthor[width=1cm,center,respectlinebreaks]\par%
 %   \vskip1.25em%
  }%
  \hbox to2cm{\hss\insertlogo\hss}
  \vskip1.25em%
  \insertverticalnavigation{2cm}%
  \vfill
  \hbox to 2cm{\hfill\usebeamerfont{subsection in
      sidebar}\strut\usebeamercolor[fg]{subsection in
      sidebar}\insertshortlecture.\insertframenumber\hskip5pt}%
  \vskip3pt%
}%

\setbeamertemplate{title page}
{
  \vbox{}
  \vskip2cm
  {\usebeamercolor[fg]{title}\usebeamerfont{title}\inserttitle\par}%
  \ifx\insertsubtitle\@empty%
  \else%
    \vskip0.25em%
    {\usebeamerfont{subtitle}\usebeamercolor[fg]{subtitle}\insertsubtitle\par}%
  \fi%     
  \vskip1em\par
%  Lecture \emph{\lecturename}\ , \insertdate\par
  \vskip0pt plus1filll
  \leftskip=0pt plus1fill\insertauthor\par
  \insertinstitute\vskip1em
}

\logo{\includegraphics[width=1cm]{njulogo.pdf}}


% Article version layout settings

\mode<article>

\makeatletter
\def\@listI{\leftmargin\leftmargini
  \parsep 0pt
  \topsep 5\p@   \@plus3\p@ \@minus5\p@
  \itemsep0pt}
\let\@listi=\@listI


\setbeamertemplate{frametitle}{\paragraph*{\insertframetitle\
    \ \small\insertframesubtitle}\ \par
}
\setbeamertemplate{frame end}{%
  \marginpar{\scriptsize\hbox to 2cm{\sffamily%
      \hfill\strut\insertshortlecture.\insertframenumber}\hrule height .2pt}}
\setlength{\marginparwidth}{1cm}
\setlength{\marginparsep}{4.5cm}

\def\@maketitle{\makechapter}

\def\makechapter{
  \newpage
  \null
  \vskip 2em%
  {%
    \parindent=0pt
    \raggedright
    \sffamily
    \vskip8pt
    {\fontsize{36pt}{36pt}\selectfont \insertshortlecture \par\vskip2pt}
    {\fontsize{24pt}{28pt}\selectfont \color{blue!50!black} \insertlecture\par\vskip4pt}
    {\Large\selectfont \color{blue!50!black} \insertsubtitle\par}
    \vskip10pt

    \normalsize\selectfont \@date\par\vskip1.5em
    \hfill Fang Yuan, ESE, NJU
  }
  \par
  \vskip 1.5em%
}

\let\origstartsection=\@startsection
\def\@startsection#1#2#3#4#5#6{%
  \origstartsection{#1}{#2}{#3}{#4}{#5}{#6\normalfont\sffamily\color{blue!50!black}\selectfont}}

\makeatother

\mode
<all>


% New useful definitions:

\newbox\mytempbox
\newdimen\mytempdimen

\newcommand\includegraphicscopyright[3][]{%
  \leavevmode\vbox{\vskip3pt\raggedright\setbox\mytempbox=\hbox{\includegraphics[#1]{#2}}%
    \mytempdimen=\wd\mytempbox\box\mytempbox\par\vskip1pt%
    \fontsize{3}{3.5}\selectfont{\color{black!25}{\vbox{\hsize=\mytempdimen#3}}}\vskip3pt%
}}

\newenvironment{colortabular}[1]{\medskip\rowcolors[]{1}{blue!20}{blue!10}\tabular{#1}\rowcolor{blue!40}}{\endtabular\medskip}

\def\equad{\leavevmode\hbox{}\quad}

\newenvironment{greencolortabular}[1]
{\medskip\rowcolors[]{1}{green!50!black!20}{green!50!black!10}%
  \tabular{#1}\rowcolor{green!50!black!40}}%
{\endtabular\medskip}



\lecture[1]{软件开发规范}{}

\begin{document}
\begin{frame}
  \maketitle
\end{frame}

\begin{frame}\frametitle<presentation>{Contents}
  \tableofcontents

\end{frame}

\section{C语言代码风格}
\begin{frame}{格式和缩进}
\begin{itemize}
    \item 代码块缩进
    \item 逻辑块之间空行
    \item 长语句断行

    \item more coding style ref. Kernel \alert{Documentation/CodingStyle}.
\end{itemize}
\end{frame}

\begin{frame}{注释}
\begin{itemize}
    \item 源文件的顶部注释: 文件名, 作者, 版本历史等
    \item 函数注释: 功能, 参数, 返回值, 错误码等
    \item 语句组上方注释
    \item 代码行右侧注释
    \item 复杂的结构体定义注释
    \item 复杂的宏定义和变量声明注释
\end{itemize}
\end{frame}

\begin{frame}{变量命名}
命名要清晰明了,易于理解
\begin{itemize}
    \item 使用完整单词、去元音缩写、首字母缩写...

        \alert{count} $\to$ \alert{cnt}, \alert{block} $\to$ \alert{blk},
        \alert{length} $\to$ \alert{len},
        \alert{internationalization} $\to$ \alert{i18n} ...
    \item 变量用下划线分割的小写字母单词, 常量用大写字母单词加下划线

        \alert{radix\_tree\_insert}, \alert{RADIX\_TREE\_MAP\_SHIFT}
    \item 全局变量和函数命名要力求详细, 宁可多用几个单词多写几个下划线

    \item 局部变量和内部函数命名可以略短,但不能太简单.
        循环变量 \alert{i}, \alert{j}, \alert{k} 等除外
    \item 不建议用汉语拼音命名标识符

\end{itemize}
    其他命名法：驼峰命名法、匈牙利命名法...
\end{frame}

\begin{frame}{函数设计原则}
\begin{itemize}

    \item 函数功能尽可能单一, 不求面面俱到 (考虑重用性和维护性)
    \item 函数内缩进层次不宜过多, 一般不超过 3 层。否则应考虑分割
    \item 函数行数一般不超过 2 屏 (48 行), 否则应考虑分割
    \item 函数名称通常应包含动词, 如
        
        \alert{get\_current()}、\alert{radix\_tree\_insert()}

    \item 重要函数上方必须加注释, 说明功能、参数、返回值、错误码等

    \item 函数内局部变量数不宜超过 10 个, 否则应考虑分割
\end{itemize}
\end{frame}

\begin{frame}{调试指令}
\begin{itemize}
    \item printf() 输出
    \item assert() 断言方法
\end{itemize}

\end{frame}


\section{软件质量}
\begin{frame}{软件质量的功能性指标}
\begin{itemize}
    \item 正确性 (Correctness)

        按照需求正确执行任务的能力
    \item 健壮性 (Robustness)
    \begin{itemize}
        \item 容错能力
        \item 恢复能力
    \end{itemize}
    \item 可靠性 (Reliability)

        平均无故障时间 (MTTF)
\end{itemize}
\end{frame}

\begin{frame}{软件质量的非功能性指标}
\begin{itemize}
    \item 性能 (Performance) $\to$ 算法复杂度分析
    \item 易用性 (Usability) (由用户判断)
    \item 清晰性 (Clarity) ，或称可读性 (Readability)
    \item 安全性 (Security) $\to$ 防止被非法入侵
    \item 可扩展性 (Extendibility) $\to$ 规模和复杂性
    \item 兼容性 (Compatibility)
    \item 可移植性 (Portability)
    \begin{itemize}
        \item 设备相关程序和设备无关程序分离
        \item 功能模块和用户界面分离
    \end{itemize}
\end{itemize}
\end{frame}

\end{document}
